\section{L1 Regularization vs Backward Elimination}
\label{sec:l1_vs_backward}

We compare two approaches for learning sparse linear mergeability predictors: (1) L1-regularized LOTO ($\lambda=1.0$) which encourages sparsity through regularization, and (2) Backward Elimination which iteratively removes metrics based on their contribution to validation performance.

\subsection{Backward Elimination}

Backward elimination is a greedy feature selection algorithm that starts with the full set of 28 metrics in the first iteration. At each iteration, we systematically remove each metric one at a time and rerun the linear optimization on the training set. We then compare the training Pearson correlation results to identify which metric, when removed, causes the least performance drop---this metric is deemed the least important. We remove that metric from the feature set and continue to the next round with the remaining metrics. Each iteration thus eliminates exactly one metric. The process terminates when removing the least important remaining metric causes a training performance drop greater than a threshold of 0.001 in Pearson correlation. Using training correlation rather than validation correlation avoids data leakage during feature selection, ensuring a fair comparison with L1 regularization which also operates solely on training data.

\subsection{Results}

Table~\ref{tab:l1_vs_backward} presents the training and validation Pearson correlation across 20 leave-one-task-out folds for each merger and feature selection approach.

\begin{table}[h]
\centering
\caption{L1 regularization vs backward elimination for mergeability prediction. We report per-fold mean and standard deviation of Pearson correlation on training and validation sets. L1 LOTO achieves higher validation performance with significantly fewer features.}
\label{tab:l1_vs_backward}
\resizebox{\textwidth}{!}{%
\begin{tabular}{@{}llccccr@{}}
\toprule
\textbf{Merger} & \textbf{Method} & \textbf{Train $r$} & \textbf{Train std} & \textbf{Val $r$} & \textbf{Val std} & \textbf{\#Features} \\
\midrule
\multirow{2}{*}{Weight Avg}
  & L1 LOTO ($\lambda$=1.0)  & 0.730 & 0.039 & \textbf{0.614} & \textbf{0.151} & 16.8 \\
  & Backward Elim.     & 0.730 & 0.038 & 0.607 & 0.128 & 26.1 \\
\midrule
\multirow{2}{*}{Arithmetic}
  & L1 LOTO ($\lambda$=1.0)  & 0.567 & 0.046 & \textbf{0.457} & \textbf{0.143} & 15.9 \\
  & Backward Elim.     & 0.467 & 0.220 & 0.334 & 0.230 & 26.6 \\
\midrule
\multirow{2}{*}{TSV}
  & L1 LOTO ($\lambda$=1.0)  & 0.753 & 0.028 & \textbf{0.657} & \textbf{0.143} & 17.1 \\
  & Backward Elim.     & 0.697 & 0.186 & 0.532 & 0.253 & 23.6 \\
\midrule
\multirow{2}{*}{TIES}
  & L1 LOTO ($\lambda$=1.0)  & 0.579 & 0.032 & \textbf{0.477} & \textbf{0.116} & 16.5 \\
  & Backward Elim.     & 0.439 & 0.104 & 0.329 & 0.166 & 26.1 \\
\midrule
\multirow{2}{*}{DARE}
  & L1 LOTO ($\lambda$=1.0)  & 0.732 & 0.030 & \textbf{0.608} & \textbf{0.189} & 15.8 \\
  & Backward Elim.     & 0.570 & 0.252 & 0.440 & 0.284 & 26.1 \\
\midrule
\midrule
\multirow{2}{*}{\textbf{Average}}
  & L1 LOTO ($\lambda$=1.0)  & \textbf{0.672} & \textbf{0.035} & \textbf{0.563} & \textbf{0.148} & \textbf{16.4} \\
  & Backward Elim.     & 0.581 & 0.160 & 0.448 & 0.212 & 25.7 \\
\bottomrule
\end{tabular}%
}
\end{table}

\subsection{Discussion}

\paragraph{L1 Regularization Outperforms Backward Elimination}
L1-regularized LOTO consistently achieves higher validation correlation across all five merging methods, with an average improvement of +0.115 over Backward Elimination (0.563 vs 0.448). This improvement is particularly pronounced for Arithmetic (+0.123), TSV (+0.125), and TIES (+0.148), where backward elimination struggles to identify the optimal feature subset.

\paragraph{Reduced Variance}
Beyond improved mean performance, L1 LOTO exhibits substantially lower variance across folds. The average validation standard deviation drops from 0.212 (Backward Elimination) to 0.148 with L1 regularization. This suggests that L1 produces more stable, generalizable coefficients that are less sensitive to which task is held out.

\paragraph{Sparsity vs. Performance Trade-off}
L1 LOTO uses only 16.4 features on average (58\% of the original 28), while Backward Elimination retains 25.7 features (92\%). Despite using 36\% fewer features, L1 achieves significantly better generalization. This demonstrates that backward elimination's greedy approach fails to identify redundant features that introduce noise, as it makes locally optimal decisions that may not be globally optimal.

\paragraph{Why Backward Elimination Underperforms}
Backward elimination suffers from several limitations: (1) it cannot recover from early mistakes---once a feature is removed, it cannot be reconsidered; (2) it evaluates features independently without accounting for feature interactions; and (3) its fixed threshold stopping criterion (0.001) may terminate too early or too late depending on the specific merger. In contrast, L1 regularization simultaneously optimizes all coefficients while encouraging sparsity, allowing correlated features to compete and the optimization to naturally identify the most predictive subset.

\paragraph{Recommendations}
Based on these results, we recommend L1-regularized LOTO over backward elimination for learning mergeability predictors. L1 regularization offers:
\begin{itemize}
    \item Higher validation correlation (+26\% relative improvement)
    \item Lower variance across folds (more reliable predictions)
    \item Sparser solutions (36\% fewer features)
    \item Simultaneous optimization avoiding greedy local optima
\end{itemize}
